\subsection{完备赋范结合代数中的指数}

在本号中，记 A 为完备赋范含幺结合代数（《一般拓扑学》，第 IX 章，§ 3，第 7 号）。于是对于 A 中的 \(\| x.y\| \leqslant \| x\|\)，有 \(\| y\|\) . \(x,y\)。

设 I 是有限集，令 \(\hat{\mathrm{P}} (\mathrm{A}^{\mathrm{I}};\mathrm{A})\) 为定义在 \(\mathbf{A}^{\mathrm{I}}\) 上、取值于 A 的带连续分量的形式幂级数的向量空间（《微分流形与解析流形》R，附录，第 5 号），其代数结构由下式给出

\[
f. g = m \circ (f, g) \quad \text { for } f, g \text { in } \hat {\mathrm{P}} (\mathrm{A} ^ {\mathrm{I}}; \mathrm{A}),
\]

其中 \(m: \mathbf{A} \times \mathbf{A} \to \mathbf{A}\) 表示 A 上的乘法。仿照第 1 号并使用 § 5 第 1 号的命题 1，定义从 \(u \mapsto \tilde{u}\) 到 \(\hat{\mathrm{A}}(\mathrm{I})\) 的含幺代数连续同态 \(\hat{\mathrm{P}}(\mathrm{A}^{\mathrm{I}}; \mathrm{A})\)，它将指标为 \(i\) 的不定元映到 \(\mathbf{pr}_{\mathrm{i}}\)；这个同态延拓了第 1 号中定义的从 \(\hat{\mathrm{L}}(\mathrm{I})\) 到 \(\hat{\mathrm{P}}(\mathrm{A}^{\mathrm{I}}; \mathrm{A})\) 的李代数同态。若 \(u = \sum_{\nu} u_{\nu}\)，其中对 \(u_{\nu} \in \mathrm{A}^{\nu}(\mathrm{I})\) 有 \(\nu \in \mathbf{N}^{\mathrm{I}}\)，则 \(\tilde{u} = \sum_{\nu} \tilde{u}_{\nu}\)，其中 \(\tilde{u}_{\nu}\) 是多项式映射 \((t_{i})_{i \in I} \mapsto u_{\nu}((t_{i}))\)。

令 \(\boldsymbol{u} = (u_{j})_{j \in I}\) 为 \(\hat{\mathbf{A}}(\mathbf{I})\) 中元素组成的有限族，令 \(v \in \hat{\mathbf{A}}(\mathbf{J})\)，并记 \(w = v \circ \boldsymbol{u}\)（\(\S 5\)，第 1 号）。则

\[
(v \circ \boldsymbol {u}) ^ {\sim} = \tilde {v} \circ \tilde {\boldsymbol {u}}.\tag{18}
\]

这是通过连续延拓 § 5 第 1 号的公式 (2)，并应用《微分流形与解析流形》R 附录第 6 号得到的。

特别地，取 \(\mathrm{I} = \{\mathrm{U}\}\)，把 A 与 \(\mathrm{A}^{(\mathrm{U})}\) 视为同一对象，并考察 \(\tilde{e}\) 中级数 \(\tilde{l}\) 和 \(e(\mathrm{U}) = \sum_{n\geqslant 1}\mathrm{U}^n /n!\) 的像 \(l(\mathrm{U}) = \sum_{n\geqslant 1}(-1)^{n - 1}\mathrm{U}^n /n\) 和 \(\hat{\mathrm{P}} (\mathrm{A};\mathrm{A})\)。于是，对 A 中的 \(\| \tilde{\mathbf{U}}^n\| \leqslant 1\)，若 \(\| x_1\dots x_n\| \leqslant \| x_1\| \dots \| x_n\|\)，则 \(x_{1},\ldots ,x_{n}\)。因此，\(\tilde{e}\)（分别为 \(\tilde{l}\)）的绝对收敛半径为无穷大（分别 \(\geqslant 1\)）。

记 \(e_{A}\)（分别为 \(l_{A}\)）为由收敛级数 \(\tilde{e}\)（分别为 \(\tilde{l}\)）定义的从 A 到 A 的解析映射（分别为从 B 到 A 的解析映射，其中 B 是 A 的开单位球），并记 \(\exp_{\mathbf{A}}(x)=1+e_{\mathbf{A}}(x)\)（对 \(x\in A\)）以及

\[
\log_ {\mathrm{A}} (x) = l _ {\mathrm{A}} (x - 1)
\]

（对 \(x \in \mathbf{A}\)，\(\| x - 1\| < 1\)）。于是

\[
\exp_ {\mathrm{A}} x = \sum_ {n \geqslant 0} \frac {x ^ {n}}{n !} (x \in \mathrm{A})\tag{19}
\]

\[
\log_ {\mathrm{A}} x = \sum_ {n \geqslant 1} (- 1) ^ {n - 1} \frac {(x - 1) ^ {n}}{n} (x \in \mathrm{A}, \| x - 1 \| <   1).\tag{20}
\]

由于 \((e\circ l)(\mathrm{U}) = (l\circ e)(\mathrm{U}) = \mathrm{U}\)（参见 § 6，第 1 号），由 (18) 有 \(\tilde{e}\circ \tilde{l} = \tilde{l}\circ \tilde{e} = \mathrm{Id}_{\mathbf{A}}\)。因此（《微分流形与解析流形》R，3.1.9）

\begin{enumerate}
\def\labelenumi{(\arabic{enumi})}
\setcounter{enumi}{20}
\tightlist
\item
\end{enumerate}

\[
\exp_ {\mathrm{A}} (\log_ {\mathrm{A}} (x)) = x \quad (x \in \mathrm{A}, \| x - 1 \| \leqslant 1)\tag{21}
\]

\[
\log_ {\mathbf {A}} (\exp_ {\mathbf {A}} (x)) = x \quad (x \in \mathrm{A}, \| x \| <   \log 2)\tag{22}
\]

因为当 \(\| x\| < \log 2\) 时有 \(\| \exp_{\mathbf{A}}(x) - 1\| \leqslant \exp \| x\| -1 < 1.\)

最后把 A 看作完备赋范李代数。于是

\[
\| [ x, y ] \| = \| x y - y x \| \leqslant 2 \| x \|. \| y \|.
\]

第 2 号的命题 1 蕴含，形式幂级数 \(\tilde{\mathbf{H}}\) 的绝对收敛域包含集合

\[
\Omega = \{(x, y) \in \mathrm{A} \times \mathrm{A} | \| x \| + \| y \| <   \frac {1}{2} \log 2 \}.
\]

因此 \(\tilde{\mathbf{H}}\) 定义了一个解析函数 \(h\colon \Omega \to \mathrm{A}\)。于是 \(h(x,y) = \sum_{r,s}\mathrm{H}_{r,s}(x,y)\)（参见 § 3，第 1 号，注 4）。

命题 3。当 \(\| x\| + \| y\| < \frac{1}{2}\log 2,\) 时，

\[
\exp_ {A} x. \exp_ {A} y = \exp_ {A} h (x, y).\tag{23}
\]

由 (18) 以及关系 \(e^{\mathrm{U}}e^{\mathrm{V}} = e^{\mathrm{H}(\mathrm{U},\mathrm{V})}\) 可知

\[
m \circ (1 + \tilde {e}, 1 + \tilde {e}) = (1 + \tilde {e}) \circ \tilde {\mathrm{H}}
\]

在 \(\hat{\mathrm{P}} (\mathrm{A}\times \mathrm{A};\mathrm{A})\) 中。因此由《微分流形与解析流形》R，3.1.9，可推出当 \((x,y)\) 充分接近 \((0,0)\) 时 (23) 成立，从而由解析延拓（《微分流形与解析流形》R，3.2.5）得到命题。

